% Options for packages loaded elsewhere
\PassOptionsToPackage{unicode}{hyperref}
\PassOptionsToPackage{hyphens}{url}
\PassOptionsToPackage{dvipsnames,svgnames,x11names}{xcolor}
\documentclass[
  12pt,
]{article}
\usepackage{xcolor}
\usepackage[a4paper, margin=2.5cm]{geometry}
\usepackage{amsmath,amssymb}
\setcounter{secnumdepth}{-\maxdimen} % remove section numbering
\usepackage{iftex}
\ifPDFTeX
  \usepackage[T1]{fontenc}
  \usepackage[utf8]{inputenc}
  \usepackage{textcomp} % provide euro and other symbols
\else % if luatex or xetex
  \usepackage{unicode-math} % this also loads fontspec
  \defaultfontfeatures{Scale=MatchLowercase}
  \defaultfontfeatures[\rmfamily]{Ligatures=TeX,Scale=1}
\fi
\usepackage{lmodern}
\ifPDFTeX\else
  % xetex/luatex font selection
  \setmainfont[]{Noto Serif}
  \setsansfont[]{Noto Sans}
  \setmonofont[]{Noto Sans Mono}
\fi
% Use upquote if available, for straight quotes in verbatim environments
\IfFileExists{upquote.sty}{\usepackage{upquote}}{}
\IfFileExists{microtype.sty}{% use microtype if available
  \usepackage[]{microtype}
  \UseMicrotypeSet[protrusion]{basicmath} % disable protrusion for tt fonts
}{}
\usepackage{setspace}
\makeatletter
\@ifundefined{KOMAClassName}{% if non-KOMA class
  \IfFileExists{parskip.sty}{%
    \usepackage{parskip}
  }{% else
    \setlength{\parindent}{0pt}
    \setlength{\parskip}{6pt plus 2pt minus 1pt}}
}{% if KOMA class
  \KOMAoptions{parskip=half}}
\makeatother
\usepackage{color}
\usepackage{fancyvrb}
\newcommand{\VerbBar}{|}
\newcommand{\VERB}{\Verb[commandchars=\\\{\}]}
\DefineVerbatimEnvironment{Highlighting}{Verbatim}{commandchars=\\\{\}}
% Add ',fontsize=\small' for more characters per line
\newenvironment{Shaded}{}{}
\newcommand{\AlertTok}[1]{\textcolor[rgb]{1.00,0.00,0.00}{\textbf{#1}}}
\newcommand{\AnnotationTok}[1]{\textcolor[rgb]{0.38,0.63,0.69}{\textbf{\textit{#1}}}}
\newcommand{\AttributeTok}[1]{\textcolor[rgb]{0.49,0.56,0.16}{#1}}
\newcommand{\BaseNTok}[1]{\textcolor[rgb]{0.25,0.63,0.44}{#1}}
\newcommand{\BuiltInTok}[1]{\textcolor[rgb]{0.00,0.50,0.00}{#1}}
\newcommand{\CharTok}[1]{\textcolor[rgb]{0.25,0.44,0.63}{#1}}
\newcommand{\CommentTok}[1]{\textcolor[rgb]{0.38,0.63,0.69}{\textit{#1}}}
\newcommand{\CommentVarTok}[1]{\textcolor[rgb]{0.38,0.63,0.69}{\textbf{\textit{#1}}}}
\newcommand{\ConstantTok}[1]{\textcolor[rgb]{0.53,0.00,0.00}{#1}}
\newcommand{\ControlFlowTok}[1]{\textcolor[rgb]{0.00,0.44,0.13}{\textbf{#1}}}
\newcommand{\DataTypeTok}[1]{\textcolor[rgb]{0.56,0.13,0.00}{#1}}
\newcommand{\DecValTok}[1]{\textcolor[rgb]{0.25,0.63,0.44}{#1}}
\newcommand{\DocumentationTok}[1]{\textcolor[rgb]{0.73,0.13,0.13}{\textit{#1}}}
\newcommand{\ErrorTok}[1]{\textcolor[rgb]{1.00,0.00,0.00}{\textbf{#1}}}
\newcommand{\ExtensionTok}[1]{#1}
\newcommand{\FloatTok}[1]{\textcolor[rgb]{0.25,0.63,0.44}{#1}}
\newcommand{\FunctionTok}[1]{\textcolor[rgb]{0.02,0.16,0.49}{#1}}
\newcommand{\ImportTok}[1]{\textcolor[rgb]{0.00,0.50,0.00}{\textbf{#1}}}
\newcommand{\InformationTok}[1]{\textcolor[rgb]{0.38,0.63,0.69}{\textbf{\textit{#1}}}}
\newcommand{\KeywordTok}[1]{\textcolor[rgb]{0.00,0.44,0.13}{\textbf{#1}}}
\newcommand{\NormalTok}[1]{#1}
\newcommand{\OperatorTok}[1]{\textcolor[rgb]{0.40,0.40,0.40}{#1}}
\newcommand{\OtherTok}[1]{\textcolor[rgb]{0.00,0.44,0.13}{#1}}
\newcommand{\PreprocessorTok}[1]{\textcolor[rgb]{0.74,0.48,0.00}{#1}}
\newcommand{\RegionMarkerTok}[1]{#1}
\newcommand{\SpecialCharTok}[1]{\textcolor[rgb]{0.25,0.44,0.63}{#1}}
\newcommand{\SpecialStringTok}[1]{\textcolor[rgb]{0.73,0.40,0.53}{#1}}
\newcommand{\StringTok}[1]{\textcolor[rgb]{0.25,0.44,0.63}{#1}}
\newcommand{\VariableTok}[1]{\textcolor[rgb]{0.10,0.09,0.49}{#1}}
\newcommand{\VerbatimStringTok}[1]{\textcolor[rgb]{0.25,0.44,0.63}{#1}}
\newcommand{\WarningTok}[1]{\textcolor[rgb]{0.38,0.63,0.69}{\textbf{\textit{#1}}}}
\setlength{\emergencystretch}{3em} % prevent overfull lines
\providecommand{\tightlist}{%
  \setlength{\itemsep}{0pt}\setlength{\parskip}{0pt}}
\usepackage{amsmath}
\usepackage{parskip}
\usepackage{fancyhdr}
\pagestyle{fancy}
\fancyhead[L]{\leftmark}
\fancyhead[R]{\thepage}
\fancyfoot{}
\usepackage{longtable,booktabs}
\usepackage{etoolbox}
\AtBeginEnvironment{longtable}{\normalsize}
\AtBeginEnvironment{verbatim}{\normalsize}
\usepackage{bookmark}
\IfFileExists{xurl.sty}{\usepackage{xurl}}{} % add URL line breaks if available
\urlstyle{same}
\hypersetup{
  pdftitle={Agents},
  pdfauthor={Mário Antunes},
  colorlinks=true,
  linkcolor={black},
  filecolor={Maroon},
  citecolor={Blue},
  urlcolor={blue},
  pdfcreator={LaTeX via pandoc}}

\title{Agents}
\usepackage{etoolbox}
\makeatletter
\providecommand{\subtitle}[1]{% add subtitle to \maketitle
  \apptocmd{\@title}{\par {\large #1 \par}}{}{}
}
\makeatother
\subtitle{Intelligent Systems II}
\author{Mário Antunes}
\date{2026}

\begin{document}
\maketitle

\setstretch{1.25}
\section{Practice Guide:
Neuro-evolution}\label{practice-guide-neuro-evolution}

In this guide, you will transition from a basic reflex agent to a
sophisticated AI capable of learning how to play Flappy Bird using
\textbf{Neuro-evolution}.

The game architecture is split into two parts:

\begin{enumerate}
\def\labelenumi{\arabic{enumi}.}
\tightlist
\item
  \textbf{Backend (Python):} Runs the game physics and logic using
  asynchronous websockets.
\item
  \textbf{Frontend (HTML5/Canvas):} Visualizes the state sent by the
  server.
\end{enumerate}

\subsection{Part 1: Setup}\label{part-1-setup}

Before writing any code, let's get the game running.

\begin{enumerate}
\def\labelenumi{\arabic{enumi}.}
\item
  \textbf{Download the Repository:} Clone the course repository:

\begin{Shaded}
\begin{Highlighting}[]
\FunctionTok{git}\NormalTok{ clone https://github.com/detiuaveiro/flappy{-}bird{-}agent.git}
\BuiltInTok{cd}\NormalTok{ flappy{-}bird{-}agent}
\end{Highlighting}
\end{Shaded}
\item
  \textbf{Create a Virtual Environment:} It is best practice to isolate
  your dependencies.

\begin{Shaded}
\begin{Highlighting}[]
\ExtensionTok{python3} \AttributeTok{{-}m}\NormalTok{ venv venv}
\BuiltInTok{source}\NormalTok{ venv/bin/activate}
\end{Highlighting}
\end{Shaded}
\item
  \textbf{Install Dependencies:} Install the required libraries.

\begin{Shaded}
\begin{Highlighting}[]
\ExtensionTok{pip}\NormalTok{ install }\AttributeTok{{-}r}\NormalTok{ requirements.txt}
\end{Highlighting}
\end{Shaded}
\item
  \textbf{Run the Game:} Start the backend server:

\begin{Shaded}
\begin{Highlighting}[]
\ExtensionTok{python} \AttributeTok{{-}m}\NormalTok{ src.backend.py }\AttributeTok{{-}{-}pipes}
\end{Highlighting}
\end{Shaded}

  Once the server is running, open the \texttt{html/play\_human.html}
  file in your web browser. Use the mouse to play the game.
\end{enumerate}

\subsection{Part 2: The Basic Agent (Reflex
Agent)}\label{part-2-the-basic-agent-reflex-agent}

Your first task is to understand the communication protocol. The server
sends the \textbf{World State} as a JSON object via websockets.

\textbf{Task:} 1. Review \texttt{play\_game} method in \texttt{play.py}.
2. Locate the message handling loop. You will see a JSON object
representing the state. 3. Implement a simple \textbf{Rule-Based Agent}.
Do not use Machine Learning yet. Write a simple \texttt{if} statement.

\textbf{Example Logic:}

\begin{Shaded}
\begin{Highlighting}[]
\CommentTok{\# A simple reflex agent}
\ControlFlowTok{if}\NormalTok{ state[}\StringTok{\textquotesingle{}bird\_y\textquotesingle{}}\NormalTok{] }\OperatorTok{\textgreater{}}\NormalTok{ state[}\StringTok{\textquotesingle{}next\_pipe\_bottom\_y\textquotesingle{}}\NormalTok{] }\OperatorTok{+} \DecValTok{50}\NormalTok{:}
\NormalTok{    action }\OperatorTok{=} \StringTok{"jump"}
\ControlFlowTok{else}\NormalTok{:}
\NormalTok{    action }\OperatorTok{=} \StringTok{"stay"}
\end{Highlighting}
\end{Shaded}

\begin{quote}
Run this agent and observe. Does it crash? Can it handle different pipe
heights?
\end{quote}

\subsection{Part 3: The Intelligent Agent
(Neuro-evolution)}\label{part-3-the-intelligent-agent-neuro-evolution}

Now we move to the core of this course: \textbf{Neuro-evolution}.

\subsubsection{1. Theoretical Background}\label{theoretical-background}

Why use Neuro-evolution for Flappy Bird?

\begin{itemize}
\tightlist
\item
  \textbf{Continuous World:} The bird's position, velocity, and pipe
  locations are continuous values (floats).
\item
  \textbf{Discrete Action:} The output is binary (Jump or Don't Jump).
\item
  \textbf{Partial Observability:} The agent doesn't need to know the
  entire map; it only needs a ``partial view'' (distance to the next
  pipes).
\end{itemize}

\textbf{The Solution: Neural Networks + Evolutionary Algorithms}

Instead of using Reinforcement Learning (like Q-Learning) which requires
a state table (hard for continuous worlds) or Gradient Descent (which
requires a dataset of ``correct'' moves), we use
\textbf{Neuro-evolution}.

\begin{itemize}
\item
  \textbf{The Brain (MLP):} We use a Multi-Layer Perceptron.
\item
  \textbf{Input:} Normalized state (Distance to pipe X, Distance to pipe
  Y, Velocity).
\item
  \textbf{Hidden Layers:} Processing units.
\item
  \textbf{Output:} A single neuron with a Sigmoid activation. If Output
  \textgreater{} 0.5 Jump.
\item
  \textbf{The Optimization (Evolution):} We do not ``train'' the network
  with backpropagation. Instead, we \textbf{evolve} the weights of the
  neural network.
\end{itemize}

\begin{enumerate}
\def\labelenumi{\arabic{enumi}.}
\tightlist
\item
  \textbf{Population:} Generate 50 birds with random weights.
\item
  \textbf{Evaluation:} Let them all play. The ``Fitness'' is the score
  (distance traveled).
\item
  \textbf{Selection:} Keep the weights of the birds that flew the
  furthest.
\item
  \textbf{Crossover \& Mutation:} Create a new generation of birds by
  mixing the weights of the best parents and adding random mutations.
\end{enumerate}

\subsubsection{2. Implementation}\label{implementation}

Open the \texttt{train.py} file. The repository is set up to use a
Python optimization library
(\href{https://pypi.org/project/pyBlindOpt/}{PyBlindOpt}; take some time
to browse the
\href{https://mariolpantunes.github.io/pyBlindOpt/pyBlindOpt.html}{documentation}).

\textbf{Your Task:} Configure the Neuro-evolution loop.

\begin{enumerate}
\def\labelenumi{\arabic{enumi}.}
\tightlist
\item
  \textbf{Define the Problem:} Map the Neural Network weights to a 1D
  vector. The optimization algorithm will try to find the vector of
  weights that maximizes the game score.
\item
  \textbf{Select an Algorithm:} The library supports several
  optimization algorithms, such as:
\end{enumerate}

\begin{itemize}
\tightlist
\item
  \textbf{DE (Differential Evolution):} Good for global search,
  maintains diversity.
\item
  \textbf{PSO (Particle Swarm Optimization):} Simulates a flock of birds
  finding food.
\item
  \textbf{GWO (Grey Wolf Optimizer):} Mimics the leadership hierarchy of
  wolves.
\item
  \textbf{CS (Cuckoo Search):} Based on the brood parasitism of cuckoo
  birds.
\end{itemize}

\subsubsection{3. Experiments}\label{experiments}

Run the training using different configurations and observe the impact
on convergence speed (how fast they learn) and stability (do they forget
how to fly?).

\textbf{Experiment A: Algorithm Comparison} Run the training for 30
generations with a population of 30 for each algorithm:

\begin{itemize}
\tightlist
\item
  Does \textbf{PSO} converge faster than \textbf{DE}?
\item
  Does \textbf{GWO} find a more stable solution?
\end{itemize}

\textbf{Experiment B: Initialization} Check the documentation for the
optimization library regarding
\href{https://mariolpantunes.github.io/pyBlindOpt/pyBlindOpt/init.html}{\emph{Initialization
Metrics}}.

\begin{itemize}
\tightlist
\item
  How does the range of initial random weights (e.g.,
  \texttt{{[}-1,\ 1{]}} vs \texttt{{[}-5,\ 5{]}}) affect the learning?
\item
  If weights start too large, the sigmoid output might saturate (always
  1 or always 0), making the bird constantly jump or never jump.
\end{itemize}

\textbf{Experiment C: Advanced Initialization Strategies} As discussed
in the lectures, standard Random Initialization can result in clusters
and gaps in the search space. The library supports advanced strategies
to mitigate this.

\begin{enumerate}
\def\labelenumi{\arabic{enumi}.}
\tightlist
\item
  \textbf{Opposition-Based Learning (OBL):}
\end{enumerate}

\begin{itemize}
\tightlist
\item
  Enable OBL initialization in your configuration.
\item
  Observation: Compare the ``Best Score'' in Generation 0 (before any
  evolution happens) between Random and OBL. Does the OBL population
  start with a ``smarter'' agent purely by checking opposite weights?
\end{itemize}

\begin{enumerate}
\def\labelenumi{\arabic{enumi}.}
\setcounter{enumi}{1}
\tightlist
\item
  \textbf{Sampling Strategies (Sobol / LHS):}
\end{enumerate}

\begin{itemize}
\tightlist
\item
  Switch the sampler to Sobol Sequence or Latin Hypercube Sampling
  (LHS).
\item
  Analysis: Does covering the search space more evenly reduce the
  variance between training runs?
\end{itemize}

\begin{enumerate}
\def\labelenumi{\arabic{enumi}.}
\setcounter{enumi}{2}
\tightlist
\item
  \textbf{OBLESA:}
\end{enumerate}

\begin{itemize}
\tightlist
\item
  Switch to the latest method named OBLESA.
\item
  Does the initial time required for OBLESA to work produces better
  results (having a bird fly the full limit with less generations)?
\end{itemize}

\subsection{Part 4: The Flappy Bird
Cup}\label{part-4-the-flappy-bird-cup}

\textbf{Friendly Competition:} Once you have explored the algorithms,
tune your hyperparameters to create an optimal agent.

\begin{itemize}
\tightlist
\item
  \textbf{Input features:} Are you feeding the agent \texttt{bird\_y}?
  Or \texttt{bird\_y\ -\ pipe\_y} (relative distance)? (Hint: Relative
  distance usually generalizes better).
\item
  \textbf{Hidden Layers:} Is a standard
  \texttt{{[}3\ inputs\ -\textgreater{}\ 5\ hidden\ -\textgreater{}\ 1\ output{]}}
  enough? Or do you need two hidden layers?
\item
  \textbf{Population Size:} Larger populations explore better but run
  slower.
\item
  \textbf{Initialization:} Consider using \textbf{OBLESA
  (Opposition-Based with Empty Space Attack)} or \textbf{QOBL
  (Quasi-Opposition)} to get a competitive edge in the early
  generations.
\end{itemize}

\textbf{Goal:} Train an agent that can reach \textbf{300 point}
consistently (fly for 5 minutes).

\end{document}
